\documentclass{article}
\usepackage[T1]{fontenc}
\usepackage{lmodern}
\usepackage{graphicx}
\usepackage{amsmath,amssymb}
\usepackage[a4paper,margin=2cm]{geometry}
\usepackage[absolute,overlay]{textpos}
\setlength{\TPHorizModule}{1mm}
\setlength{\TPVertModule}{1mm}
\usepackage[indonesian]{babel}
\usepackage{hyperref}
\setlength{\emergencystretch}{2em}
% unit: Halaman 19

\begin{document}

% === Halaman 19 (llm) ===

\section*{Kelebihan Mode $ECB$}

\begin{enumerate}
    \item Karena tiap blok plainteks dienkripsi secara independen, maka kita tidak perlu mengenkripsi pesan secara sekuensial/linier.
\end{enumerate}

\begin{itemize}
    \item Kita dapat mengenkripsi 5 blok pertama, kemudian blok-blok di akhir, dan kembali ke blok-blok di tengah dan seterusnya.
    \item Mode $ECB$ cocok untuk mengenkripsi isi arsip (\textit{file}) yang diakses secara acak, misalnya arsip-arsip basisdata.
    \item Jika basisdata dienkripsi dengan mode $ECB$, maka sembarang \textit{record} dapat dienkripsi secara independen dari \textit{record} lainnya (dengan asumsi setiap \textit{record} terdiri dari sejumlah blok diskrit yang sama banyaknya).
\end{itemize}

\end{document}
